% ISOLATED PROVISIONAL SECTION. Not included by manuscript347.tex.
% Source346A frozen in PR287; separate audit pending at preparation.
% Do not incorporate as accepted before the parent supplies its actual verdict.
\section{Provisional support-size exclusion from source346}
This section reproduces a frozen candidate under separate analytical audit. It is isolated from all accepted conclusions above and is not a premise for them. Its inclusion in a final manuscript requires the actual source346 verdict and a subsequent manuscript audit. It uses the actual $hPh$ relation and requires no marginal matching.

\begin{theorem}[Provisional source346 statement]\label{support346}
Let $0<\rho<1$, $h>0$, $\E hPh=1$, $\pi=hPh\geq b>0$ and $\E\pi^s\leq2$, where $s=(1+\rho)/(1-\rho)$. Write
\[
 B=2^{1/s},\quad a_0=(1+\rho)/2,\quad d=a_0^L,\quad
 \theta=1-1/(2s),\quad \kappa=2^{1/s}a_0(1+\rho^3).
\]
For every nonempty $A$ of cardinality $m$ with $md<1$, put
\[
 R_L(m)=\frac{\sqrt B}{b}\,
 \frac{m^\theta\kappa^{L/2}}{(1-md)^2}.
\]
Then for every $x$, $K(x,A)\leq R_L(m)/(1+R_L(m))$. The killed kernel $K_A$ has spectral radius at most this bound, and
\[
 \frac{J(A)}{\nu(A)}\geq\frac1{1+R_L(m)}.
\]
If $\kappa<1$ and $0<\tau<\min\{-\log a_0,-\log\kappa/(2\theta)\}$, these estimates give uniformly $R_L(m)\to0$ for $m\leq e^{\tau L}$. In the explicit range $0<\rho\leq1/4$ and $m\leq e^{\rho L/1920}$,
\[
 R_L(m)\leq\frac{9\sqrt B}{b}\,e^{-\rho L/640}.
\]
\end{theorem}
\begin{proof}
The maximal entry and every diagonal entry of $P$ equal $d$. Thus $\pi(y)\geq d h(y)^2$, and, with $\alpha_0=2^{-L}$,
\begin{align*}
 I_x(A)&:=\sum_{y\in A}P(x,y)h(y)
 \leq\sqrt d\sum_{y\in A}\sqrt{\pi(y)}\\
 &\leq\sqrt{Bd}\,m^{1-1/(2s)}\alpha_0^{-1/(2s)}.
\end{align*}
The last inequality is finite-sum H\"older and $\sum_y\pi(y)^s\leq2/\alpha_0$. This upper bound alone cannot control the normalized transition, since its denominator may also be small.

For the outside contribution $O_x(A)=\sum_{y\notin A}P(x,y)h(y)$, use $h\geq b/g$. Since $P(x,A)\leq md$, Cauchy--Schwarz with weights $P(x,y)$ on $A^c$ gives
\[
 (1-md)^2\leq
 \left(\sum_{y\notin A}\frac{P(x,y)}{g(y)}\right)
 \left(\sum_{y\notin A}P(x,y)g(y)\right)
 \leq\left(\sum_{y\notin A}\frac{P(x,y)}{g(y)}\right)(P^2h)(x).
\]
Therefore $O_x(A)\geq b(1-md)^2/(P^2h)(x)$. The semigroup gives $P^2h=P_{\rho^{3/2}}Qh$. Point evaluation of $P_{\rho^{3/2}}$ on uniform $L^2$ has squared norm
\[
 \frac{P_{\rho^3}(x,x)}{\alpha_0}=(1+\rho^3)^L.
\]
Because $\norm{Qh}_2=1$, we obtain $(P^2h)(x)\leq(1+\rho^3)^{L/2}$. Dividing the upper bound on $I_x$ by the lower bound on $O_x$ yields $I_x/O_x\leq R_L(m)$, exactly with the stated $\kappa$. Now $K(x,A)=I_x/(I_x+O_x)$, proving the pointwise bound.

Every row sum of $K_A$ is at most $R/(1+R)$, so its spectral radius is at most that number. Summing $K(x,A^c)\geq1/(1+R)$ over $x\in A$ against $\nu$ proves the flux bound. Equivalently, reversibility gives the same upper Rayleigh bound for functions supported on $A$ in $L^2(\nu)$.

If $m\leq e^{\tau L}$ in the specified range, then $md\to0$ and $m^\theta\kappa^{L/2}\to0$, uniformly. For the explicit subrange,
\begin{align*}
 \log\kappa
 &=-\frac{2\rho}{1+\rho}\log2+\log(1+\rho)+\log(1+\rho^3)\\
 &\leq\rho\left(\frac{17}{16}-\frac85\log2\right)
 <-\rho/240,
\end{align*}
using $\log2>2/3$. As $\theta<1$, $m^\theta\kappa^{L/2}\leq e^{-\rho L/640}$. Moreover
\[
 md\leq[(5/8)e^{1/7680}]^L\leq(2/3)^L,
\]
because $\log(16/15)\geq1/16>1/7680$. Thus $(1-md)^{-2}\leq9$, completing the estimate.
\end{proof}

At $\kappa=1$, the formula still gives a positive uniform escape bound for singletons, at least $[1+\sqrt B/(b(1-a_0)^2)]^{-1}$. For $\kappa>1$ the criterion can become asymptotically uninformative even for singletons: at $\rho=1/2$, $\kappa=27\,2^{1/3}/32>1$ because $39366>32768$. This is a limitation of the estimate, not an admissible collapsing example. Small stationary mass need not imply small cardinality, so even an accepted version of this section would not close the original gap question or transfer a critical moment bound to $r$.
